% ============================================================
% OP3_1_ordering_recon_v1.tex  (PROPOSAL-ONLY RECONNAISSANCE)
% OP3.1 ordering mechanism: map, interface ledger, programme
% formalisation. NO preprint edits, NO mechanism claims, NO
% promise of a fast derivation.
% ============================================================
\section*{OP3.1 ordering-mechanism reconnaissance (proposal-only)}

\subsection*{0. Scope honesty}
OP3.1 (``the ordering mechanism itself, including the microscopic
derivation of the closure functions from bare~P3'') is the deepest
layer of the transition block.
No fast derivation is promised and none is attempted here.
The realistic deliverables, in the session's established mould,
are: (A)~a \emph{necessary interface ledger} -- conditions any
ordering mechanism must satisfy, assembled from already-recorded
Level~A/B structures on both sides of the gap; and (B)~a
\emph{programme formalisation} of the configuration-space
measure/saddle route (zero-mode-D1 style), with the
stochastic-quantisation route recorded and fenced.
OP3.1 remains Open throughout; at most ``Open; interface
requirements recorded / programme sharpened''.

\subsection*{1. Existing anchors (both sides of the gap)}
\emph{Below the gap (input side):} bare~P3 (condensate postulate);
the pre-ordered configuration regime Proposition (Level~A,
\S pre\_lorentzian): no effective Lorentzian metric, no causal
cone, no emergent temporal direction prior to ordering -- so the
mechanism may not presuppose clock-time evolution.
\emph{Apparatus:} the recorded extremal-action / saddle-dominance
principle (Level~A, \S qs\_extremal\_action): Euclidean weight
$\mathcal W[q]\sim\exp(-S_E[q]/S_0)$
(eq:qs\_extremal\_weight\_E), saddle dominance for
$|\Delta S_E|/S_0\gg1$.
\emph{Above the gap (output side):} the ordered-branch Euclidean
EFT functional (eq:ordered\_branch\_euclidean\_eft) with data
$Z_u$, $\mathcal C_n$, $W(u)$ (and corrections); the
closure-function mapping (\S closure\_function\_mapping):
$F=e^{\beta_\phi\phi}$, $\omega$ from the amplitude stiffness,
$U(\phi)=W(u_\infty e^{\beta_\phi\phi})$ up to recorded
normalisation freedoms.
\emph{Downstream consistency:} OP3.2 T-reversal lemma~A,
no-bounce sublemma, conditional persistence criterion; OP3.3
bookkeeping ledger ($\rho_{\rm onset}$ as ordered-branch effective
density); the no-phantom bound (eq:w\_phi).

\subsection*{2. Task A0 (audit; prerequisite)}
Locate and record the precise form in which the \emph{bare}~P3
functional/measure is currently specified in the article.
If only the ordered-branch EFT is pinned while the bare functional
is specified at postulate level, then the first genuine
sub-problem is OP3.1a: \emph{state the admissible bare functional
class} (fields, symmetries, locality, positivity) from which the
derivation is to proceed.
Without A0 any ``derivation'' claim is ill-posed.

\subsection*{3. Kernel A: ordering-mechanism interface ledger
(M1--M6; necessary conditions, no mechanism)}
\emph{(M1) Arena and admissible inputs.}
Inputs are bare~P3 data on Euclidean configurations only.
By the pre-ordered-regime Proposition, no emergent-time or
causal-cone structure may be presupposed; any ``relaxation''
parameter must be a formal internal device, never the emergent
physical time.
\emph{(M2) Apparatus compatibility.}
A measure/saddle mechanism must be expressible in the recorded
$S_0$-weighted Euclidean form and must reduce to the recorded
saddle-dominance statement in the regime $|\Delta S_E|/S_0\gg1$;
``selection of the ordered branch'' then means \emph{measure
concentration} on ordered configurations, a claim requiring proof,
not assertion.
\emph{(M3) Output interface (success criterion).}
Closing OP3.1 means computing from the bare functional:
(a)~existence of ordered extrema/minimisers;
(b)~their dominance over disordered configurations;
(c)~the ordered-branch EFT data ($Z_u$, $\mathcal C_n$, $W(u)$,
leading corrections) by expansion around the dominant
configuration -- which, via the recorded closure-function mapping,
is equivalent to deriving $F,\omega,U$ up to the recorded
freedoms.
Postulating $W(u)$ is not closing OP3.1.
\emph{(M4) Symmetry and orientation interface.}
The derived dynamics must exhibit the OP3.2 time-reversal
covariance at equation level; branch-orientation selection must
enter through boundary/measure data rather than hidden bias, and
must either satisfy or explicitly delegate the assumptions of the
recorded persistence criterion.
\emph{(M5) Cosmological hand-off.}
The mechanism must furnish $\rho_{\rm onset}$ as an ordered-branch
effective density within the FLRW reduction, in the sense fixed by
the OP3.3 ledger; it must make no energy statements at the
pre-emergent Euclidean level.
\emph{(M6) Consistency floors.}
The induced closure functions must respect the recorded
positivity/no-ghost expectations of the scalar--tensor closure and
the no-phantom bound (eq:w\_phi) on the ordered branch.

\subsection*{4. Kernel B: saddle-selection programme
(Definition / Open-problem block; all Open)}
\emph{Working requirements (not axioms):} a specified bare
functional class (A0); the recorded $S_0$-weight; a notion of
ordered vs disordered configuration sectors (order parameter
$u,n_A$ at coarse level).
\emph{Definition (ordering-selection problem).}
Within the recorded weight, establish (i)~existence of ordered
saddles/minimisers of the bare functional on the constrained
configuration space; (ii)~dominance/measure concentration of the
ordered sector; (iii)~extraction of the EFT data of M3(c) by
controlled expansion.
\emph{Conditional existence scaffold (to be activated only after
A0).}
If the bare functional is bounded below and coercive on the
constrained space, the direct method yields existence of
minimisers; the sharp mathematical comparison is then between
$O(4)$-symmetric and ordering ($O(4)\to O(3)$) minimisers.
Stated as a conditional scaffold; the hypotheses are exactly what
A0 must make checkable.
\emph{Failure modes (recorded in advance):}
no concentration regime ($|\Delta S_E|/S_0\not\gg1$ where it
matters); flat/degenerate directions spoiling dominance;
boundary-data dependence of the selected sector (consistent with
the article's boundary-value framing); order-parameter
coarse-graining ambiguities.
\emph{Firewall.}
Nothing in this programme presupposes a time variable; ``selection''
is a statement about the Euclidean measure, not a process in time.

\subsection*{5. Route C (recorded, deferred): stochastic-quantisation
/ relaxational formalisms}
Recorded as in the article (``more speculative'').
Explicit circularity firewall for any future use: the fictitious
relaxation parameter of stochastic quantisation is a formal device
for representing the same Euclidean measure; it must not be
identified with, or used to generate, the emergent physical time,
on pain of circularity.
Admissible role: computational reformulation of Kernel~B only.
Deferred; no literature inserted before source verification
(Parisi--Wu and a standard review would be the candidates).

\subsection*{6. What is NOT proposed}
No ordering mechanism; no derivation of $U,\omega,F$; no claim
that the measure selects the ordered branch; no change to
OP3.1/OP3.2/OP3.3 statuses; no time-language in the pre-ordered
regime; no promise of a timeline.

\subsection*{7. Risk register}
Forbidden: ``ECT derives the ordering mechanism / the Big Bang'';
``the measure selects the ordered saddle'' (without proof);
clock-time or relaxation-in-time language for the pre-ordered
regime; treating the EFT functional as if it were the bare one.
Safe: ``necessary interface conditions, recorded'';
``programme sharpened; mechanism open'';
``conditional scaffold, hypotheses pending A0''.

\subsection*{8. Questions for GPT}
(Q1)~Approve the M1--M6 interface ledger as the article-facing
deliverable (Level~B conditional bookkeeping, AN3/OP3.3 mould)?
(Q2)~Approve Kernel~B as a zero-mode-D1-style
Definition/Open-problem block, including the conditional existence
scaffold in its ``activated only after A0'' form -- or keep the
scaffold proposal-only until A0 is done?
(Q3)~Task A0: run it as the immediate next mini-round (audit of
how bare~P3 is currently pinned in the text), before any
insertion?
(Q4)~Placement: both blocks adjacent to the OP3.1 item of the
decomposition (post-transition subsection), with one-line echoes
only; nothing inside the quantum-sector saddle subsection?
(Q5)~Status target ``Open; interface requirements recorded /
programme sharpened'' -- agreed?
